\documentclass[11pt,a4paper]{article}
\usepackage[utf8]{inputenc}
\usepackage[T1]{fontenc}
\usepackage{lmodern}
\usepackage{geometry}
\geometry{margin=1in}
\usepackage{hyperref}
\usepackage{booktabs}
\usepackage{graphicx}
\usepackage{microtype}
\usepackage{setspace}
\onehalfspacing

\title{Future Weapons and Artificial Intelligence: Autonomous Warfare, Strategic Stability, and the Governance Crisis}
\author{}
\date{}

\begin{document}
\maketitle

\begin{abstract}
The integration of artificial intelligence into military systems is no longer speculative. Fully autonomous loitering munitions have been operational since 1989, AI-driven targeting pipelines are fielded by the Pentagon, and drone swarms are reshaping battlefields in Ukraine. Yet the international governance framework remains paralyzed: the UN Convention on Certain Conventional Weapons (CCW) operates by consensus, allowing any single state to block progress, while the major military powers---the United States, China, and Russia---pursue aggressive autonomous weapons programs while resisting binding legal constraints. This paper surveys the current landscape of AI-enabled military systems, examines the divergence between technological deployment and diplomatic regulation, analyzes the strategic stability implications of autonomous warfare, and identifies the accountability gap that emerges when lethal decisions are delegated to algorithms. Drawing on empirical case studies from the Israeli Harpy loitering munition to Ukraine drone warfare and China's intelligentized warfare doctrine, we argue that autonomous weapons create a governance trilemma: states cannot simultaneously achieve military effectiveness, legal accountability under international humanitarian law, and strategic stability.
\end{abstract}

\section{Introduction}

In 1989, Israel Aerospace Industries fielded the Harpy, a fully autonomous loitering munition that searches for radar emissions and destroys them without human approval. In 2024, the U.S. Department of Defense requested 1.8 billion dollars for the second tranche of its Replicator Initiative, an ambitious program to field thousands of attritable autonomous systems across all domains. In the same year, China's People's Liberation Army published doctrinal texts describing autonomous uncrewed platforms as the preferred type of battlefield equipment in the evolution toward intelligentized warfare. These are not isolated developments. They represent the leading edge of a transformation in military affairs that may prove as consequential as the advent of nuclear weapons or mechanized warfare.

The scholarly and policy debate has lagged behind the technology. Much of the literature on lethal autonomous weapons systems (LAWS) remains either speculative---projecting futures of artificial general intelligence commanding armies---or narrowly technical, focused on specific subsystems. What is missing is a grounded assessment that integrates three distinct literatures: the technical literature on AI reliability and failure modes; the strategic studies literature on arms races and crisis stability; and the international law literature on accountability and compliance. Existing surveys tend to treat these dimensions separately. This paper addresses that gap by analyzing them as interconnected elements of a single governance trilemma.

\section{Conceptual Framework: What Counts as Autonomy?}

A persistent obstacle to both analysis and governance is definitional ambiguity. The term autonomous weapon is used to describe systems ranging from terminal defense systems like Phalanx---which autonomously intercepts incoming anti-ship missiles within a narrow spatial envelope---to hypothetical future systems employing artificial general intelligence for strategic planning.

\subsection{Three Levels of Human Involvement}

A more productive framework distinguishes systems by the nature and timing of human involvement in the targeting cycle:

\begin{itemize}
\item \textbf{Human-in-the-loop}: The system identifies or tracks targets, but a human operator must approve each specific engagement.
\item \textbf{Human-on-the-loop}: The system can select and engage targets autonomously within pre-defined parameters, but a human operator can intervene to override.
\item \textbf{Human-out-of-the-loop}: The system searches for, identifies, and engages targets without any human intervention or override capability after launch.
\end{itemize}

\subsection{AI-Enabled Decision Support vs. Lethal Autonomy}

A second critical distinction separates lethal autonomous weapons from AI-enabled command, control, and decision support systems.

\section{The Current Landscape: Systems, Programs, and Deployment}

\subsection{The United States: Speed, Scale, and Ambiguity}

The U.S. Department of Defense has made AI and autonomous systems central to its modernization strategy. Project Maven demonstrated that machine learning could dramatically reduce the time required to analyze drone and satellite imagery for potential targets. The Replicator Initiative aims to field thousands of attritable autonomous systems across multiple domains.

\subsection{China: Doctrine, Civil-Military Fusion, and the Race for Intelligentized Warfare}

China's approach to military AI differs fundamentally from the United States in its doctrinal ambition, state-directed integration, and ideological framing. It is essential to distinguish between doctrinal aspiration and operational capability: Western assessments consistently note significant gaps between Chinese ambition and demonstrated fielded capability.

\subsection{Russia and Ukraine: The Crucible of Robotic Warfare}

The Russia-Ukraine conflict has become a crucible for autonomous and semi-autonomous systems. Russia has deployed Lancet and KUB loitering munitions extensively across eastern Ukraine.

\subsection{Israel: Operational Excellence and Export Leadership}

Israel stands at the cutting edge of operational autonomous weapons deployment. The IAI Harpy, which entered service in the early 1990s following development that began in the late 1980s, is the first fully autonomous loitering munition.

\subsection{Europe and Other States: Regulatory Ambition and Operational Caution}

The European Union and key European states occupy a distinct position in the autonomous weapons landscape: advanced technological capability paired with stronger regulatory ambition than the major military powers.

\subsection{Summary of Major Autonomous Weapons Programs}

\begin{table}[h]
\centering
\begin{tabular}{lllll}
\toprule
Country & System & Autonomy Level & Domain & Status \
\midrule
United States & Project Maven & Human-in-the-loop & Air/Space & Operational \
United States & Replicator Initiative & Mixed & Multi-domain & In development \
United States & XQ-58A Valkyrie & Human-on-the-loop & Air & Demonstration \
China & Wing Loong II & Human-on-the-loop & Air & Operational \
Russia & Lancet/KUB & Human-on-the-loop & Air & Combat deployed \
Russia & Marker UGV & Human-on-the-loop & Ground & Testing \
Israel & Harpy & Human-out-of-the-loop & Air & Operational \
Israel & Harop & Human-on-the-loop & Air & Operational \
Israel & Iron Dome & Human-out-of-the-loop & Air & Operational \
Turkey & Kargu-2 & Human-on-the-loop & Air & Combat deployed \
Turkey & Bayraktar TB2 & Human-in-the-loop & Air & Combat deployed \
UK & Autonomous Warrior & Human-in-the-loop & Ground/Sea & Testing \
France/Germany & FCAS/MGCS & Human-in-the-loop & Air/Ground & Development \
\bottomrule
\end{tabular}
\end{table}

\section{The Governance Crisis: Diplomacy Lags Behind Technology}

The gap between technological deployment and diplomatic regulation is widening. Discussions on LAWS under the UN Convention on Certain Conventional Weapons began in 2013. The CCW operates by consensus, meaning any single state can block progress.

\subsection{The Verification Problem}

Any governance framework for autonomous weapons faces a fundamental verification challenge. Unlike nuclear weapons, which produce detectable signatures, autonomous capabilities are primarily software-de
